\documentclass[11pt,letterpaper]{article}
\usepackage[left=.87in,right=.87in,top=.77in,bottom=.77in]{geometry}
\usepackage[T1]{fontenc}
\usepackage{lmodern,amsmath,amssymb,amsthm,booktabs,microtype,xcolor}
\usepackage[colorlinks=true,allcolors=teal!65!black]{hyperref}
\newtheorem{theorem}{Theorem}
\newtheorem{lemma}[theorem]{Lemma}
\newtheorem{corollary}[theorem]{Corollary}
\newcommand{\Q}{\mathbb Q}\newcommand{\Z}{\mathbb Z}\newcommand{\F}{\mathbb F}
\newcommand{\ord}{v_p}\newcommand{\ds}{s_p}
\newcommand{\src}[2]{\href{https://github.com/VladimirReshetnikov/ProveIt/blob/cc919c585f90be8e0fba7a1b90c70afdea39998e/Analysis/GeneralizedRiemannHypothesis/#1}{#2}}
\linespread{0.968}\title{\vspace{-1.4em}Sharp sparse root counts beyond the first Frobenius jump}
\author{A research supplement to the ProveIt project}
\date{27 November 2026}
\begin{document}\maketitle
\begin{abstract}
How many roots of a polynomial with at most $t$ monomials can lie strictly inside the $p$-adic exponential disk? A first jump to $p^2$ occurs at $t=p+2$. We prove that the exact maximum through $2p+1$ terms is $p^2+\lfloor(t-p-2)/2\rfloor$. For arbitrary $t$ we give a finite determinant formula for the entire depth profile on each fixed support. Its input requires only the exponent polynomial modulo an explicit power of $p$. We prove that the necessary number of residue digits attains our upper bound for infinitely many term budgets at each prime. The proofs use integral moment recurrences and the stabilization of Frobenius on a finite exponent algebra. All root counts are geometric and the results are unconditional.
\end{abstract}

\section{The extremal problem}
Normalize $\ord(p)=1$. In $\overline{\Q}_p$ set
\[
 D=\{x:\ord(x-1)>1/(p-1)\},\qquad
 D_r=\{x:\ord(x-1)\ge r\},\qquad r>1/(p-1).
\]
For a finite set $E$ of distinct nonnegative integers, let $C_E(r)$ and $C_E^\circ$ be the maxima of the number of roots in $D_r$ and $D$, counted with multiplicity, of a nonzero $P(X)=\sum_{n\in E}c_nX^n$. Coefficients may belong to any finite extension of $\Q_p$. A nonzero center $A_0$ gives the same counts on the disks defined by $x/A_0$ in place of $x$. This follows by scaling the coefficients. Write $B_p(t)=\sup_{|E|\le t} C_E^\circ$.

The repository's \src{Article/sections/smd.tex}{singular-moduli section} bounds local root counts by factorial-weighted moments. Its \src{Research/sparse-padic-roots-and-graph-towers/fewnomial/paper.tex}{open-disk supplement} proves a triangular-number law through $p+1$ terms and $B_p(p+2)=p^2$. It asks how to find the sharp bound for larger supports, where Frobenius retains nilpotent information, and how to compute the fixed-support profile at every depth. A \src{Research/padic-support-rank/report.tex}{closed-disk supplement} gives a residue rank test for $t\le p$. We use the same analytic reduction and solve the fixed-support question for all $t$.

\begin{theorem}\label{thm:stairs}
For every prime $p$ and every $p+2\le t\le 2p+1$,
\[
 \boxed{\displaystyle B_p(t)=p^2+\left\lfloor\frac{t-p-2}{2}\right\rfloor.}
\]
Each maximum is realized by a polynomial with integer coefficients and simple roots.
\end{theorem}
At $p=5$ the exact maxima for seven through eleven terms are $25,25,26,26,27$. The formula gives $121,121,122,122,123,123,124,124,125,125,126$ at $p=11$. Once the first jump has occurred, two further terms buy one additional root up to the stated endpoint.

\section{Moment coordinates and a comparison principle}
For $t=|E|$ put $H_E(T)=\prod_{n\in E}(T-n)$, $\mathcal A_E=\Z_p[T]/H_E(T)$, with power basis $1,\ldots,T^{t-1}$. Over any coefficient field the moments $M_k=\sum_{n\in E} c_n n^k$ are values $L(T^k)$ for a linear functional $L$ on the exponent algebra. Every functional occurs by Vandermonde inversion. Scale $L$ so its power-basis values are integral and at least one is a unit. Then all moments are integral. Write $s(k)=\ds(k)$ for the sum of the base-$p$ digits and $r_0=1/(p-1)$. For rational $r>r_0$ put $\delta=r-r_0$ and choose $b$ with $\ord(b)=r$. The analytic coordinate $X=\exp(bY)$ takes the closed unit disk bijectively onto $D_r$, with inverse $\log(X)/b$ and nonzero derivative. The coefficient of $Y^k$ is $b^k M_k/k!$. By Legendre's formula its valuation is
\begin{equation}\label{eq:weights}
 \ord(M_k)+w_r(k),\qquad w_r(k)=\delta k+s(k)/(p-1).
\end{equation}
This tends to infinity with $k$, with the convention $v_p(0)=+\infty$. Weierstrass preparation gives the geometric root count as the largest minimizing index in \eqref{eq:weights}.

\begin{lemma}\label{lem:compare}
Set $d=\max_{0\le j<t}s(j)$ and $a=\max\{1,\lceil(d-1)/(p-1)\rceil\}$. A minimizing index $k\ge t$ satisfies $s(k)<d$ and $\ord(M_k)<a$. If $T^k$ modulo $p^a\mathcal A_E$ belongs to the span over $\Z/p^a\Z$ of powers $T^\ell$ with $w_r(\ell)<w_r(k)$, this $k$ cannot minimize.
\end{lemma}
\begin{proof}
Compare with a unit moment at some $j<t$. Then
\[
 \ord(M_k)\le w_r(j)-w_r(k)
 \le \delta(j-k)+(d-s(k))/(p-1).
\]
Use integrality, $s(k)\ge1$, and $k>j$. For the second assertion, lift a relation to the integral exponent algebra. Its error evaluates to an element of valuation at least $a$. One term on the right must therefore have moment valuation at most $\ord(M_k)$, contradicting minimality.
\end{proof}

\section{One large residue class}
\begin{proof}[Upper bound in Theorem~\ref{thm:stairs}]
It suffices to handle supports of size $t$ since zero coefficients permit padding. When all exponents occupy one residue, factor $P(X)=X^uQ(X^{p^e})$ after removing the largest common $p$-power in exponent differences. The power map is injective on $D$ by the logarithm and preserves multiplicity. The resulting support has at least two distinct residues. Thus it suffices to prove the open bound for such supports, by bounding every rational closed subdisk of $D$.

For $t\le2p+1$, $d\le p$, so Lemma~\ref{lem:compare} applies with $a=1$. A minimizing $k\ge t$ has digit sum $\le p-1$. Work in $R=\mathcal A_E/p\mathcal A_E$, a product of factors $\F_p[T]/(T-c)^{m_c}$. If $m_c\le p$ throughout, then $T^p=c$ on each factor. Writing $i$ for the units digit of $k$ and $m$ for the sum of its higher digits, $T^k=T^{i+pm}$ there. The latter index is $<p^2$ since $i+m\le p-1$. When $k\ge p^2$ this lowers the weight strictly, by \eqref{eq:weights}. This excludes those indices.

Otherwise exactly one residue has multiplicity $p+u>p$. Multiply $P$ by a monomial to place that class at zero. Put $b_0=t-p-u$, so $1\le b_0\le p$, $p+u\le2p$. Write $R=R_0\times R_*$ where $R_0=\F_p[T]/T^{p+u}$ and $R_*$ has dimension $b_0$ with $T$ a unit. Every $T^{p^j}$ for $j\ge1$ acts by the same componentwise scalar on $R_*$, since its factors have multiplicity $\le p$. On $R_0$ every $T^{p^j}$ for $j\ge2$ vanishes, and $T^{2p}=0$. If $k\ge p^2$ and $m\ge2$ then $T^k=T^{i+pm}$ in both factors. This again gives a strictly lower weight. For $m=1$ replace $k=i+p^j$, $j\ge2$, by $i+p^2$. A reduction is strict unless $j=2$.

It remains to handle $k=p^2+i$. If $i\ge u$, then $T^{p^2+i}=T^{p+i}$, so $k$ is excluded. If $i\ge b_0$, express $T^i$ on $R_*$ as a combination of $1,T,\ldots,T^{b_0-1}$. Multiplication by $T^{p^2}$ kills $R_0$, and gives a global relation for $T^k$ in powers $T^{p^2+j}$ with $j<i$. Each has strictly smaller weight. Thus surviving high indices satisfy
\[
 i\le\min(u,b_0)-1\le \lfloor(t-p-2)/2\rfloor.
\]
The initial indices $k<t$ obey the same overall bound since $t-1\le2p\le p^2$. Taking $r\downarrow r_0$ captures all roots in $D$ for each polynomial.
\end{proof}

\begin{proof}[Sharpness in Theorem~\ref{thm:stairs}]
Put $i=\lfloor(t-p-2)/2\rfloor$, $s=p+i+1$, $b_0=i+1$, $C=p^2+i$. Take $s$ exponents congruent to zero and $b_0$ congruent to one. They require $p+2i+2\le t$ slots. Then
\[
 R=\F_p[T]/T^s\ \times\ \F_p[y]/y^{i+1},\qquad T=1+y\quad\text{on the second factor}.
\]
On the second factor take the coefficient of $y^i$. On the first assign to $T^k$ for $0\le k<s$ the negative of $\binom{k}{i}$ (zero when $k<i$). The resulting functional $\lambda$ has
\[
 \lambda(T^k)=0\ (k<s),\qquad
 \lambda(T^k)=\binom{k}{i}\bmod p\ (k\ge s).
\]
Since $i<p$, Lucas's theorem gives $\lambda(T^C)=1$. For $k\ge s$ write $j=k\bmod p$. A nonzero $\lambda(T^k)$ needs $j\ge i$. If $s(k)\le i$, this would force $k=i<s$. At digit sum $i+1$, the only possibilities with $j\ge i$ are $k=i+1$ or $k=i+p^h$ for $h\ge1$. Those with $k<C$ are smaller than $s$ or equal $p+i=s-1$. Thus the functional annihilates every earlier index in the lexicographic order $(s(k),k)$.

Lift $\lambda$ to an integral functional on the full power basis, choosing $M_0=0$ exactly. Integer lifts give rational coefficients for $P$ by Vandermonde inversion. With $r>r_0$ rational sufficiently close to $r_0$ one has
\[
 w_r(C)<\min\{(i+2)/(p-1),\ 1+1/(p-1)\},
\]
since $i+1\le p-1$ including $p=2$. Any index of smaller weight is annihilated modulo $p$, by the preceding order argument. Its coefficient has valuation at least $1+1/(p-1)$ for a positive index, and the constant is zero. All other moments are integral and $M_C$ is a unit. This proves there are at least $C$ geometric roots. The upper bound forces equality.
\end{proof}

To obtain simple roots and full support in this construction, perturb its rational coefficients by numbers of sufficiently high $p$-adic valuation. All moment perturbations are uniformly small, hence the last minimizing index at the chosen rational $r$ is preserved. After making the allowed coefficients nonzero and removing a common monomial, vary the constant term to avoid the finitely many critical values of the nonconstant polynomial. Then clear denominators. Our sharpness supports may be chosen to contain zero, so the resulting polynomial is squarefree.

\section{All term budgets and all depths}
The next theorem reduces arbitrary exponent sizes and an infinite sequence of moments to explicit finite data.

\begin{theorem}\label{thm:finite}
For a support of size $t\ge2$ keep $d,a$ from Lemma~\ref{lem:compare}, and set
\[
 h=\lceil\log_p t\rceil,\quad q=h+a-1,\quad
 K=\max\{t-1,(d-1)p^q\}.
\]
The full profile $C_E(r)$ and $C_E^\circ$ depends only on $H_E(T)\bmod p^a$. To compute it, write $v_k$ for the column of $T^k$ in the power basis of $\mathcal A_E$. For each $t$-subset $J\subseteq\{0,\ldots,K\}$ with nonzero $\Delta_J=\det(v_j)_{j\in J}$ put
\[
 F_J(r)=\ord(\Delta_J)+\sum_{j\in J}\bigl(rj-\ord(j!)\bigr).
\]
Then
\[
 \boxed{\displaystyle C_E(r)=\max\{\max J:F_J(r)=\min_I F_I(r)\}.}
\]
For $C_E^\circ$ minimize the pair $(F_J(r_0),\sum_{j\in J}j)$ lexicographically and take the largest index among all minimizing sets. Integer coefficients on the fixed support attain each maximum, with all nonzero roots simple.
\end{theorem}

\begin{proof}
First, reductions of $T^{p^j}$ stabilize modulo $p^a$ for $j\ge q$. Indeed $T^{p^h}$ is constant at its residue on every factor of $\mathcal A_E/p$ since every factor multiplicity is $\le t\le p^h$. Thus $T^{p^{h+1}}\equiv T^{p^h}\pmod p$. In the integral algebra, $x\equiv y\pmod{p^m}$ implies $x^p\equiv y^p\pmod{p^{m+1}}$ for $m\ge1$, by the binomial expansion. Iterating proves stabilization. Replace each summand $p^j$ in the digit expansion of $k$, for $j>q$, by $p^q$. This gives a congruent power $T^{k'}$, with $k'\le s(k)p^q$, $s(k')\le s(k)$ after carries. At a candidate $k>K$, Lemma~\ref{lem:compare} gives $s(k)\le d-1$, whence $k'<k$ and $w_r(k')<w_r(k)$. Such $k$ cannot minimize.

For precision, if $H_E\equiv H_{E'}\pmod{p^a}$ then all their power-remainder columns match modulo $p^a$. Give a functional on the second algebra exactly the same first $t$ normalized moments as on the first. Every moment then agrees modulo $p^a$ in the valuation ring of the coefficient field. Each moment of valuation below $a$ keeps its valuation. Every index $k\ge t$ with moment valuation $\ge a$ exceeds a coefficient valuation from the first $t$ indices, by the strict comparison in Lemma~\ref{lem:compare}. The minimizing indices match in both directions. This works even if some coefficient valuations are fractional.

For the determinant formula fix rational $r>r_0$ and scale column $v_k$ to $z_k=b^kv_k/k!$, as in \eqref{eq:weights}. Over a finite extension containing $b$, $F_J(r)$ is the valuation of the $J$-minor of this finite matrix. The columns have full rank because $v_0,\ldots,v_{t-1}$ form a basis. Let $I$ minimize $F$. By Cramer's rule the basis $z_I$ expresses every $z_k$ with integral coefficients, so $z_I$ is an integral basis of the lattice spanned by all these columns. A column has nonzero image modulo the maximal ideal of this lattice exactly when it lies in some minimizing basis. In fact a nonzero reduction extends to a basis among the spanning reduced columns, and a determinant is a unit precisely when its reductions form a basis.

Scale any nonzero functional so its minimum valuation on the columns is zero. Columns with zero reduction cannot attain this minimum. Conversely, for a minimizing basis containing a column $z_k$, take the coordinate functional at that column. Its minimum is attained at $k$. Because no index above $K$ can minimize, this proves the formula for the largest possible index. The coordinate functional is a scalar multiple of a rational one: annihilate the other $t-1$ rational vectors $v_j$ in that basis. This gives rational polynomial coefficients by Vandermonde inversion. Perturbation and clearing denominators as above give full support and simple nonzero roots.

The $F_J$ are finitely many affine functions with rational coefficients. For irrational $r$, the minimizing sets are constant in a neighborhood of $r$ (identical functions may tie throughout). Monotonicity and the rational case give the value there. A single witness from a slightly deeper rational disk attains it. Immediately to the right of $r_0$, minimization uses the stated pair. For each polynomial the count on $D$ is already attained on some rational closed subdisk. Conversely a witness for any such closed disk supplies as many roots in $D$. These facts identify the open-disk maximum.
\end{proof}

The dependence assertion in Theorem~\ref{thm:finite} permits any convenient integer representatives of an exponent multiset modulo $p^a$, choosing representatives distinct within a repeated class. The determinant computation uses ordinary integer remainders for those representatives; individual determinant valuations need not be determined modulo $p^a$. For a prescribed residue polynomial that is known to arise, one can also enumerate residue multisets to find a representative. Maximizing the resulting profiles over the finitely many multisets of $t$ elements of $\Z/p^a\Z$ computes the unrestricted maximum at every depth. Extra slots permit zero coefficients, so this includes all smaller supports. The lower envelope gives every transition point. This finite algorithm gives decidability without asserting practical running time for large $t$.

\begin{corollary}[A uniform degree bound]\label{cor:degree}
At every term budget $t\ge2$ and every prime,
\[
 B_p(t)\le(t-1)p^a < p\,t(t-1).
\]
\end{corollary}
\begin{proof}
Put $m=p^a$. If every exponent has the same residue modulo $m$, take representatives $c,c+m,\ldots,c+(t-1)m$ and remove $X^c$. Otherwise take the first available nonnegative representatives in each class. Each multiplicity is then at most $t-1$, so the resulting degree is at most $(t-1)m-1$. The profile of these supports equals the original profile by Theorem~\ref{thm:finite}. For the last inequality use $d\le h(p-1)$, hence $a\le h$, and $p^h<pt$.
\end{proof}
There is also a direct computational shortcut: use the canonical representatives just described, with minimum exponent shifted to zero. Expand $X^n=(1+Y)^n$. The same minimum-minor argument applies to the finite columns $(\binom{n}{k})_n$, using the weight $rk$ instead of $rk-\ord(k!)$. This computes the full profile with at most $1+(t-1)p^a$ columns from ordinary binomial coefficients. The argument is simply the Newton polygon on $\ord(Y)\ge r$.

\section{The precision bound is genuinely needed}
\begin{theorem}\label{thm:optimal}
For each prime, the power $p^a$ of Theorem~\ref{thm:finite} cannot in general be reduced to $p^{a-1}$, for infinitely many $t$. This obstruction already occurs for the open-disk maximum.
\end{theorem}
\begin{proof}
Take $t=p^h$, with $h\ge2$ for odd primes and $h\ge3$ for $p=2$. Here $d=h(p-1)$, so $a=h$ for odd $p$ and $a=h-1$ for $p=2$. Compare
\[
 E=\{0,1,\ldots,t-1\},\qquad
 E'=\{0,1,\ldots,t-2,N\},\qquad N=t-1+p^{a-1}.
\]
Their exponent multisets agree modulo $p^{a-1}$. The first support has maximum precisely $t-1$, by degree and $(X-1)^{t-1}$. Subtract from $X^N$ its Taylor polynomial at $1$ through degree $t-2$. This yields an $E'$-polynomial
\[
 (X-1)^{t-1}Q(X-1),\quad
 Q(0)=\binom N{t-1},\quad Q'(0)=\binom Nt .
\]
The first binomial has valuation $h-a+1$ by the carries in adding $p^{a-1}$ to $p^h-1$. The second is a unit by Lucas's theorem. Hensel's lemma gives a root $y$ of $Q$ of positive valuation, with $\ord(y)=h-a+1$ by comparing its constant and linear terms. This is $1$ for odd $p$ and $2$ for $p=2$, strictly above $r_0$. The count on $E'$ is at least $t$.
\end{proof}

For $p=2$, the eight-term supports $\{0,\ldots,7\}$ and $\{0,\ldots,6,9\}$ yield exactly seven and eight roots. Indeed the latter Taylor remainder is $Y^7(36+9Y+Y^2)$ with $Y=X-1$. All nine roots of a degree-nine polynomial on that support cannot lie in $D$: their sum is zero by the missing $X^8$ coefficient, whereas nine elements reducing to $1$ sum to $1$ in residue characteristic two.

\section{Reproducible checks and scope}
The supplied standard-library Python program reconstructs rational witnesses via Lagrange interpolation, clears denominators and perturbs the constant. It counts their open-disk roots directly from $P(1+Y)$, using the smallest index minimizing $\ord([Y^k]P(1+Y))+k/(p-1)$. The smallest index gives the right-hand limit at the boundary. For each of 32 witnesses at $p=2,3,5,7,11,13,17$ it separately certifies global simplicity by squarefree reduction at another prime. A second calculation scans residue jets for all 4,002 multisets at $p=2,3,5$ in the staircase range. A third compares 100 truncated moment determinant calculations with determinants of the entire finite binomial expansion.

The general exponential-polynomial bounds of van der Poorten and Rumely give $(n-2)p$ for order $3\le n<p$ over $\Q_p$ and treat more general frequencies~\cite{vdr}. Avenda\~no and Krick study how many roots belong to a specified local field~\cite{ak}. Our geometric counts on a single multiplicative disk allow arbitrary finite extensions for the roots. The staircase and optimal residue precision answer local questions in the cited ProveIt reports. This argument leaves the general sharp local-field fewnomial problem open. Neither the zeta function nor a hypothesis about its zeros enters the proof. The literature comparison establishes the mathematical scope of the extension. Full publication priority has not been established.

Beyond $2p+1$ slots, two different residue classes can both retain a nilpotent after one Frobenius application. What is the next closed formula for $B_p(t)$? For budgets that are not powers of $p$, which admit a precision smaller than that in Theorem~\ref{thm:finite}? The finite test supplies a way to generate and certify examples for both questions.

\begin{thebibliography}{9}
\bibitem{vdr} A. J. van der Poorten and R. S. Rumely,
\emph{Zeros of $p$-adic exponential polynomials II},
J. London Math. Soc. (2) 36 (1987), 1--15.
\href{https://doi.org/10.1112/jlms/s2-36.1.1}{doi:10.1112/jlms/s2-36.1.1}.
\bibitem{ak} M. Avenda\~no and T. Krick,
\emph{Sharp bounds for the number of roots of univariate fewnomials},
J. Number Theory 131 (2011), 1209--1228.
\href{https://doi.org/10.1016/j.jnt.2011.01.006}{doi:10.1016/j.jnt.2011.01.006}.
\end{thebibliography}
\end{document}
